\documentclass[12pt]{article}
\usepackage[utf8]{inputenc}
\usepackage[T1]{fontenc}
\usepackage{hyperref}
\title{Zarphatic Language Resources for K-12 Curriculum Development}
\author{Compiled by Claude Sonnet 5.5 (claude-sonnet-5-5) for Andr\'as Kornai}
\date{}
\begin{document}
\maketitle

\section{Introduction}
Zarphatic (ISO 639-3 zrp; Glottocode zarp1238), also called Judeo-French, is an extinct Jewish language of medieval northern France and parts of western Germany. It is attested from the 11th century, notably in glosses to the Torah and Talmud by Rashi and Moshe HaDarshan, and it died out around the 14th century. It was written in Hebrew script, and many linguists regard it as essentially Old French in a distinct writing system. The name was coined by Solomon Birnbaum from the Hebrew \emph{Tzarfat}, meaning France. There are no native speakers, so any K-12 use would be heritage and revitalization material built from historical sources.

\section{Primary Educational Resources}
No school curriculum or primers for Zarphatic were found in the search results.

\section{Linguistic and Reference Materials}
Wikipedia gives an overview of the language: \url{https://en.wikipedia.org/wiki/Zarphatic_language}. The Jewish Language Project has a Judeo-French page: \url{https://www.jewishlanguages.org/judeo-french}. The Rosetta Panglossia entry lists the language under the code zrp: \url{https://rosettapanglossia.longnow.org/index.php?title=Zrp}. Verbix lists a page for Judeo-French: \url{https://docs.verbix.com/Languages/Judeo-French}.

\section{Religious and Community Resources}
The principal attestations are the medieval Hebrew-script glosses on biblical and Talmudic texts, but the searches found no online edition. The Jewish Language Project page above is the nearest community-facing resource.

\section{Digital and Audio Resources}
No audio exists, since the language is extinct. A Wiktionary-derived word list is at \url{https://kaikki.org/dictionary/All languages combined/meaning/Z/Za/Zarphatic.html}. It covers the word ``Zarphatic'' only, not the language's vocabulary.

\section{Curriculum Development Context}
Material for a heritage unit would come from the medieval glosses and the scholarly literature. That would suit a unit within Jewish history, Hebrew or French language classes, not a full-language programme. Nothing in the search results shows that such a programme exists.
\end{document}